\documentclass[11pt,a4paper]{article}

\usepackage[margin=1in]{geometry}
\usepackage{booktabs}
\usepackage{longtable}
\usepackage{array}
\usepackage{multirow}
\usepackage{xcolor}
\usepackage{colortbl}
\usepackage{tikz}
\usepackage{enumitem}
\usepackage{amsmath}
\usepackage{hyperref}
\usepackage{float}
\usepackage{caption}
\usepackage{tabularx}

\usetikzlibrary{arrows.meta, positioning, shapes.geometric, calc}

\definecolor{yes}{HTML}{2E7D32}
\definecolor{no}{HTML}{C62828}
\definecolor{partial}{HTML}{F57F17}
\definecolor{degraded}{HTML}{E65100}
\definecolor{lightgray}{HTML}{F5F5F5}
\definecolor{sectionblue}{HTML}{1565C0}

\newcommand{\yes}{\textcolor{yes}{\textbf{Yes}}}
\newcommand{\no}{\textcolor{no}{\textbf{No}}}
\newcommand{\prt}{\textcolor{partial}{\textbf{Partial}}}
\newcommand{\dgr}{\textcolor{degraded}{\textbf{Degraded}}}
\newcommand{\na}{\textcolor{gray}{N/A}}

\newcolumntype{C}[1]{>{\centering\arraybackslash}m{#1}}
\newcolumntype{L}[1]{>{\raggedright\arraybackslash}m{#1}}

\title{\textbf{Social Grassroots Networking: Transport Connectivity Analysis}\\[0.5em]
\large Address Types, Peer Topologies, and Mobility Scenarios}
\author{Social Grassroots Networking --- Transport Layer Documentation}
\date{\today}

\begin{document}

\maketitle
\tableofcontents
\newpage

% =============================================================================
\section{Introduction}
% =============================================================================

This document provides a comprehensive analysis of connectivity between Social Grassroots Networking (SGN) peers across all combinations of network address types, peer topologies, and mobility scenarios. SGN uses two transport layers:

\begin{itemize}
  \item \textbf{BLE (Bluetooth Low Energy)} --- short-range, no Internet required, used for local peer discovery and communication
  \item \textbf{LibP2P} (UDX/TCP) --- used for global connectivity beyond BLE range
\end{itemize}

In addition to direct transports, SGN supports \textbf{application-level peer relay}: a connected peer C can forward end-to-end encrypted messages between peers A and B when no direct path exists.

\subsection{Key Constraints}

\begin{enumerate}
  \item \textbf{No store-and-forward}: If a peer is unreachable (even via relay), the message fails. No caching, no queuing. Relay is real-time only.
  \item \textbf{Single-hop peer relay}: A connected peer C \textbf{may} forward messages from A to B if C can reach B. Messages are end-to-end encrypted; C cannot read them. Maximum one relay hop (no multi-hop chains A$\to$C$\to$D$\to$B).
  \item \textbf{Transport priority}: BLE first, then libp2p (direct), then peer relay (last resort).
\end{enumerate}

\subsection{Address Types Under Analysis}

\begin{table}[H]
\centering
\begin{tabular}{llp{7cm}}
\toprule
\textbf{Address Type} & \textbf{Abbreviation} & \textbf{Description} \\
\midrule
BLE only & BLE & Bluetooth Low Energy; $\sim$10--100m range, no IP stack \\
Private IPv4 & PIv4 & Behind NAT (e.g., \texttt{192.168.x.x}, \texttt{10.x.x.x}); not routable from Internet \\
Public IPv4 & PuIv4 & Directly routable IPv4; rare for mobile devices \\
Link-local IPv6 & LLv6 & \texttt{fe80::/10}; only valid on local link, not routable \\
Public IPv6 & PuIv6 & Globally routable IPv6; increasingly common on mobile carriers \\
\bottomrule
\end{tabular}
\caption{Network address types evaluated in this document.}
\end{table}

\subsection{Peer Topology Definitions}

We evaluate four peer topologies:

\begin{enumerate}[label=\textbf{T\arabic*}]
  \item \textbf{Friend Next to Me} --- Both peers are physically co-located ($<$30m). Both BLE and Internet transport are potentially available.
  \item \textbf{Friend Far Away} --- Peers are geographically separated. Only Internet transport is possible. No BLE.
  \item \textbf{Relay Peer (``Friend in the Middle'')} --- A third peer C exists between peers A and B. We evaluate whether C can facilitate A$\leftrightarrow$B communication via transport-level relay (Circuit Relay) or application-level message forwarding.
  \item \textbf{Server} --- Either (a) public relay/bootstrap infrastructure, or (b) a custom always-on supernode with a stable public IP.
\end{enumerate}

% =============================================================================
\section{Transport Mechanics}
% =============================================================================

\subsection{BLE Transport}

\begin{itemize}
  \item \textbf{Range}: Typically 10--30m indoors, up to 100m outdoors with line of sight.
  \item \textbf{Discovery}: Each device advertises a unique service UUID derived from its Ed25519 public key. Scanning discovers all nearby BLE devices; service discovery confirms SGN peers.
  \item \textbf{Data}: After GATT service discovery and ANNOUNCE exchange, data flows over the SGN GATT characteristic.
  \item \textbf{No IP required}: BLE operates below the IP layer. Works without any Internet or WiFi connectivity.
  \item \textbf{Symmetric}: Both peers must be in range; there is no multi-hop BLE.
\end{itemize}

\subsection{LibP2P Transport}

\begin{itemize}
  \item \textbf{Transports}: Two transport protocols are available:
    \begin{itemize}
      \item \textbf{UDX} (UDP-based, custom protocol) --- primary transport, preferred for performance
      \item \textbf{TCP} --- secondary, used for IPFS ecosystem interoperability and as firewall fallback
    \end{itemize}
  \item \textbf{Security}: Noise protocol for all connections (both UDX and TCP).
  \item \textbf{Identity}: LibP2P uses its own Ed25519 keypair to generate a \texttt{PeerId}, \textbf{separate} from the SGN application-layer Ed25519 public key. This requires maintaining a double mapping:
    \begin{enumerate}
      \item BLE Device ID $\leftrightarrow$ LibP2P PeerId
      \item LibP2P PeerId $\leftrightarrow$ SGN Public Key
    \end{enumerate}
  \item \textbf{Addressing}: Multiaddr format (e.g., \texttt{/ip4/1.2.3.4/tcp/4001/p2p/12D3Koo...})
  \item \textbf{Application protocol}: \texttt{/gever/1.0.0} --- each message sent via a new stream (not persistent connection)
  \item \textbf{DHT}: Supported but \textbf{disabled} (\texttt{enableDht = false}). SGN never uses Kademlia for peer discovery.
  \item \textbf{mDNS}: \textbf{Enabled} (\texttt{enableMdns = true}). Provides automatic local network peer discovery without BLE.
\end{itemize}

\subsubsection{LibP2P NAT Traversal Stack}

LibP2P provides a multi-layered NAT traversal strategy:

\begin{enumerate}
  \item \textbf{AutoNAT v2}: Probes reachability by asking other peers to dial back. Classifies the local NAT type (public, restricted, symmetric).
  \item \textbf{Circuit Relay v2}: Public peers volunteer as relays. NATed peers reserve relay slots. Traffic flows: A $\to$ Relay Peer $\to$ B. \textit{Unlike dedicated relay servers, these are libp2p peers themselves.}
  \item \textbf{AutoRelay}: Automatically discovers and connects to available relay peers when AutoNAT detects the node is behind NAT.
  \item \textbf{DCUtR (Direct Connection Upgrade through Relay)}: After establishing a relayed connection, peers coordinate hole-punching to upgrade to a direct connection.
\end{enumerate}

\textbf{Address priority hierarchy} (highest to lowest):
\begin{enumerate}
  \item IPv6/UDX --- best performance, globally routable
  \item IPv4/UDX --- common, good performance
  \item IPv4/TCP --- IPFS interop, wider firewall compatibility
  \item Circuit Relay --- last resort, highest latency
\end{enumerate}

\subsubsection{LibP2P Bootstrap Requirement}

LibP2P \textbf{requires bootstrap peers} for relay discovery:

\begin{verbatim}
/ip4/40.160.9.115/tcp/4001/p2p/12D3KooWKnDdG3...
/ip6/2604:2dc0:101:100::138f/tcp/4001/p2p/12D3KooWKnDdG3...
\end{verbatim}

These bootstrap peers serve as:
\begin{itemize}
  \item Initial connection points for newly started nodes
  \item Relay reservation targets (via Circuit Relay v2)
  \item AutoNAT probing partners
\end{itemize}

Without the bootstrap peer being reachable, a NATed node cannot discover relays and is limited to direct connections or mDNS.

\subsubsection{LibP2P Peer Discovery}

LibP2P does \textbf{not} discover peers autonomously (DHT is disabled). Discovery happens via:
\begin{enumerate}
  \item \textbf{BLE ANNOUNCE packets} --- peers discovered via BLE include their libp2p multiaddrs
  \item \textbf{mDNS} --- automatic discovery on the same LAN (no BLE required)
  \item \textbf{Out-of-band exchange} --- manual address sharing (QR code, link)
  \item After discovery, the app calls \texttt{connectToHost(hostId, hostAddrs)} with learned addresses
\end{enumerate}

\subsection{Application-Level Peer Relay}

In addition to the direct transports (BLE and libp2p), SGN supports \textbf{application-level message forwarding} through connected peers.

\begin{itemize}
  \item \textbf{How it works}: If Peer A cannot reach Peer B directly (neither BLE nor libp2p), but Peer C is connected to both A and B, then A can send the message to C, and C forwards it to B.
  \item \textbf{Cross-transport}: C can receive via BLE from A and forward via libp2p to B (or any combination). This bridges transports that could not otherwise interoperate.
  \item \textbf{End-to-end encrypted}: C sees only the encrypted envelope. It cannot read, modify, or inspect message content.
  \item \textbf{Single hop only}: Only A$\to$C$\to$B is permitted. No multi-hop chains (A$\to$C$\to$D$\to$B).
  \item \textbf{Real-time only}: If C cannot reach B right now, the forward fails immediately. C does not queue or cache messages.
  \item \textbf{Any peer can relay}: No special server role required. Any connected peer can forward messages for others.
\end{itemize}

This is distinct from libp2p's Circuit Relay v2, which operates at the transport level. Application-level relay works \textit{across} transport boundaries and does not require C to have a public IP.

% =============================================================================
\section{Connectivity: Friend Next to Me (T1)}
% =============================================================================

Both peers are within BLE range ($<$30m). Each row represents the peer's network address type. We evaluate whether each transport works.

\begin{table}[H]
\centering
\small
\begin{tabular}{L{2.2cm} C{1.5cm} C{1.3cm} C{1.3cm} C{4.5cm}}
\toprule
\textbf{My Addr} & \textbf{Peer Addr} & \textbf{BLE} & \textbf{LP2P} & \textbf{Notes} \\
\midrule
BLE only & BLE only & \yes & \no & BLE works. No IP = no libp2p. \\
\rowcolor{lightgray}
BLE only & PIv4 & \yes & \no & I have no IP; cannot use libp2p. \\
BLE only & PuIv4 & \yes & \no & Same---I have no IP stack. \\
\rowcolor{lightgray}
BLE only & LLv6 & \yes & \no & Same. \\
BLE only & PuIv6 & \yes & \no & Same. \\
\midrule
PIv4 & PIv4 & \yes & \yes$^*$ & \textbf{mDNS discovers peer on same LAN} without relay. If different LANs: needs bootstrap $\to$ relay. \\
\rowcolor{lightgray}
PIv4 & PuIv4 & \yes & \yes & Direct TCP or UDX to public peer. No relay needed. \\
PIv4 & LLv6 & \yes & \prt & mDNS might discover on same link. UDX over link-local unreliable. \\
\rowcolor{lightgray}
PIv4 & PuIv6 & \yes & \dgr & Cross-protocol. Needs relay peer that bridges IPv4$\leftrightarrow$IPv6. Not guaranteed. \\
\midrule
PuIv4 & PuIv4 & \yes & \yes & Direct. No relay needed. \\
\rowcolor{lightgray}
PuIv4 & PIv4 & \yes & \yes & Peer connects to me directly or via relay. \\
PuIv4 & LLv6 & \yes & \no & Link-local not routable. mDNS possible on same link only. \\
\rowcolor{lightgray}
PuIv4 & PuIv6 & \yes & \dgr & Cross-protocol relay uncertain. Relay peer must be dual-stack. \\
\midrule
LLv6 & LLv6 & \yes & \prt & mDNS on same link. UDX \textit{might} bind link-local. \\
\rowcolor{lightgray}
LLv6 & Any public & \yes & \no & Cannot reach bootstrap or peer. \\
\midrule
PuIv6 & PuIv6 & \yes & \yes & Direct IPv6/UDX. Optimal. \\
\rowcolor{lightgray}
PuIv6 & PIv4 & \yes & \dgr & Cross-protocol via relay. Relay must be dual-stack. \\
PuIv6 & PuIv4 & \yes & \dgr & Same cross-protocol issue. \\
\bottomrule
\end{tabular}
\caption{Co-located connectivity (T1). $^*$mDNS provides same-LAN discovery without relay.}
\end{table}

\subsection{T1 Analysis}

\textbf{Key insight}: When peers are co-located, \textbf{BLE is always available} regardless of network configuration. BLE provides the reliable connectivity floor.

LibP2P adds value beyond BLE when co-located:
\begin{itemize}
  \item Higher throughput than BLE (UDX/TCP vs GATT characteristic writes)
  \item \textbf{mDNS} discovers peers on the same LAN without needing BLE at all---useful when devices are on the same WiFi but beyond BLE range
  \item TCP fallback works through firewalls that block UDP
\end{itemize}

\textbf{Cross-protocol weakness}: When one peer has only IPv4 and the other only IPv6, libp2p requires a dual-stack relay peer. Since relay peers are other libp2p nodes (not dedicated infrastructure), dual-stack availability is not guaranteed.

\textbf{Peer relay impact on T1}: Minimal. Co-located peers already have BLE as a reliable direct transport. Peer relay is unnecessary when direct connectivity exists.

% =============================================================================
\section{Connectivity: Friend Far Away (T2)}
% =============================================================================

Peers are geographically separated. BLE is not available. Only libp2p can work.

\begin{table}[H]
\centering
\small
\begin{tabular}{L{2.2cm} C{1.8cm} C{1.3cm} C{6cm}}
\toprule
\textbf{My Addr} & \textbf{Peer Addr} & \textbf{LP2P} & \textbf{Notes} \\
\midrule
BLE only & Any & \no & No IP = no libp2p. \textbf{Unreachable.} \\
\rowcolor{lightgray}
PIv4 & BLE only & \no & Peer has no IP. \textbf{Unreachable.} \\
PIv4 & PIv4 & \dgr & Both NATed. Need bootstrap peer reachable. AutoRelay finds relay. DCUtR attempts hole-punch. More moving parts than dedicated relay. \\
\rowcolor{lightgray}
PIv4 & PuIv4 & \yes & Direct to peer. Can use TCP or UDX. \\
PIv4 & LLv6 & \no & Not routable. \\
\rowcolor{lightgray}
PIv4 & PuIv6 & \dgr & Need dual-stack relay peer. Not guaranteed in libp2p relay pool. \\
\midrule
PuIv4 & PIv4 & \yes & Peer uses relay then DCUtR to me. \\
\rowcolor{lightgray}
PuIv4 & PuIv4 & \yes & Direct. Both UDX and TCP work. \\
PuIv4 & PuIv6 & \dgr & Cross-protocol. Relay or direct only if both dual-stack. \\
\midrule
\rowcolor{lightgray}
LLv6 & Any & \no & Cannot reach bootstrap or any remote peer. \textbf{Unreachable.} \\
\midrule
PuIv6 & PIv4 & \dgr & Cross-protocol relay required. \\
\rowcolor{lightgray}
PuIv6 & PuIv4 & \dgr & Same issue. Direct only if peer also has IPv6. \\
PuIv6 & PuIv6 & \yes & Direct IPv6/UDX. Best case. \\
\bottomrule
\end{tabular}
\caption{Remote connectivity (T2). Cross-protocol scenarios limited by relay peer heterogeneity.}
\end{table}

\subsection{T2 Analysis}

\textbf{Critical dependency}: Without BLE, the \textit{only} direct transport is libp2p. If either peer lacks a routable IP address (BLE-only or link-local IPv6), communication is \textbf{impossible} unless a peer relay can bridge the gap.

\textbf{Peer relay impact on T2}: Significant for BLE-only peers. If Peer A is BLE-only but has a co-located Peer C who has both BLE and libp2p connectivity, C can forward messages from A to the remote Peer B. This transforms previously ``unreachable'' scenarios into reachable ones---provided a suitable relay peer C exists.

\subsubsection{NAT Traversal Success Rates}
\label{sec:nat}

Hole-punching success (via DCUtR) depends on NAT types:

\begin{table}[H]
\centering
\begin{tabular}{lcccc}
\toprule
& \textbf{Full Cone} & \textbf{Restricted} & \textbf{Port Restricted} & \textbf{Symmetric} \\
\midrule
\textbf{Full Cone} & \yes & \yes & \yes & \yes \\
\textbf{Restricted} & \yes & \yes & \yes & \prt \\
\textbf{Port Restricted} & \yes & \yes & \prt & \no \\
\textbf{Symmetric} & \yes & \prt & \no & \no \\
\bottomrule
\end{tabular}
\caption{Hole-punching success by NAT type combination. Symmetric NAT (common on carrier-grade NAT / mobile networks) is the hardest case.}
\end{table}

When DCUtR hole-punching fails, libp2p falls back to Circuit Relay. Relay always works as long as both peers can reach the relay peer. However:
\begin{itemize}
  \item \textbf{Relay availability}: Depends on the bootstrap peer being online. If bootstrap is down, NATed peers have no relay.
  \item \textbf{Symmetric NAT $\leftrightarrow$ Symmetric NAT}: Relay-only; no direct upgrade possible. Higher latency ($\sim$50--150ms added round-trip depending on relay location).
  \item \textbf{Single relay}: In practice, SGN uses primarily one bootstrap peer as relay---a single point of failure.
\end{itemize}

% =============================================================================
\section{Connectivity: Relay Peer (T3) --- ``Friend in the Middle''}
\label{sec:relay-peer}
% =============================================================================

\textbf{Scenario}: Peer A wants to reach Peer B. Peer C is a third peer connected to both A and B. Can C help?

SGN supports relay at \textbf{two levels}:
\begin{enumerate}
  \item \textbf{Transport-level} (Circuit Relay v2): LibP2P's built-in relay. Requires C to have a public IP. Operates within the libp2p transport only.
  \item \textbf{Application-level} (Peer Forwarding): SGN's own forwarding. Any peer C can relay. Works across transport boundaries (BLE$\to$libp2p, etc.). End-to-end encrypted.
\end{enumerate}

\begin{table}[H]
\centering
\small
\begin{tabular}{L{4cm} C{1.2cm} C{1.2cm} C{5cm}}
\toprule
\textbf{Scenario} & \textbf{Circuit Relay} & \textbf{App Relay} & \textbf{Details} \\
\midrule
C is public, A \& B NATed, all on libp2p & \yes & \yes & Circuit Relay: transport-level, Noise-encrypted. App relay: also works, end-to-end encrypted. \\
\rowcolor{lightgray}
C is NATed too, all on libp2p & \no & \yes & Circuit Relay requires public IP. \textbf{App relay works regardless}---C just needs libp2p connections to both A and B. \\
C has BLE to A, libp2p to B & \no & \yes & \textbf{Cross-transport forwarding}: C bridges BLE and libp2p. A$\to$C via BLE, C$\to$B via libp2p. Previously impossible. \\
\rowcolor{lightgray}
C is the bootstrap peer & \yes & \yes & Bootstrap peer serves as both transport relay and potential app relay. \\
A$\leftrightarrow$C (BLE), C$\leftrightarrow$B (BLE), B out of A's range & \no & \yes & \textbf{BLE range extension}: C forwards between two BLE-only peers. End-to-end encrypted. \\
\rowcolor{lightgray}
A$\to$C$\to$D$\to$B (multi-hop) & \no & \no & Multi-hop forbidden. Single relay hop only. \\
C connected to A but B is offline & \no & \no & Real-time only. No queuing at C. \\
\bottomrule
\end{tabular}
\caption{Relay peer analysis. Transport-level (Circuit Relay) vs application-level (peer forwarding).}
\end{table}

\subsection{T3 Analysis}

Application-level peer relay significantly expands the relay peer's utility compared to Circuit Relay alone:

\begin{itemize}
  \item \textbf{Cross-transport bridging}: The biggest gain. A BLE-only peer can now communicate with a remote libp2p peer through a co-located intermediary. This was previously impossible---BLE and libp2p were completely separate worlds.
  \item \textbf{No public IP requirement}: Unlike Circuit Relay, any connected peer can forward, including NATed peers. This means relay capacity grows organically with the network.
  \item \textbf{BLE range extension}: Three peers in a line (A--C--B) where C is in BLE range of both A and B can now communicate. Effective BLE range doubles to $\sim$60m indoors / $\sim$200m outdoors.
  \item \textbf{End-to-end encrypted}: C sees only encrypted envelopes. No trust required in the relay peer beyond connectivity.
\end{itemize}

\textbf{Limitations}:
\begin{itemize}
  \item \textbf{Single hop}: Only one relay peer in the path. Cannot chain relays.
  \item \textbf{Real-time}: If C cannot reach B immediately, the forward fails. No store-and-forward.
  \item \textbf{Relay peer must exist}: Requires a Peer C that happens to be connected to both A and B. In sparse networks, this may not exist.
\end{itemize}

% =============================================================================
\section{Server Scenarios (T4)}
% =============================================================================

\subsection{T4a: Bootstrap Peer (LibP2P Infrastructure)}

The bootstrap peer is a server at \texttt{40.160.9.115} with both IPv4 and IPv6 addresses.

\begin{table}[H]
\centering
\small
\begin{tabular}{L{3.5cm} C{1.5cm} C{6cm}}
\toprule
\textbf{Role} & \textbf{Works?} & \textbf{Details} \\
\midrule
Initial peer discovery & \yes & Newly started nodes connect to bootstrap to join the network. \\
\rowcolor{lightgray}
Relay server (Circuit Relay) & \yes & NATed peers reserve relay slots on bootstrap. Primary relay path. \\
AutoNAT probe target & \yes & Bootstrap dials back to test reachability. \\
\rowcolor{lightgray}
Single point of failure & \textcolor{no}{\textbf{Risk}} & If bootstrap goes down: no relay for NATed peers, no AutoNAT, new peers cannot join. \\
\bottomrule
\end{tabular}
\caption{Bootstrap peer roles.}
\end{table}

\textbf{Single point of failure}: The bootstrap peer is a critical dependency. If it goes offline:
\begin{itemize}
  \item Already-connected peers with direct connections continue working
  \item NATed peers lose their relay path (no new relay reservations)
  \item New peers cannot join the network at all
  \item AutoNAT cannot classify NAT type
\end{itemize}

\subsection{T4b: Custom Always-On Supernode}

A server with a stable public IP running a full libp2p peer.

\begin{table}[H]
\centering
\small
\begin{tabular}{L{3.5cm} C{1.5cm} C{6.5cm}}
\toprule
\textbf{Capability} & \textbf{Works?} & \textbf{Details} \\
\midrule
Transport relay & \yes & Can serve as Circuit Relay v2. Any peer can reserve relay slots. \\
\rowcolor{lightgray}
Always-on peer & \yes & Reachable by any peer with IP connectivity. \\
Bootstrap redundancy & \yes & Can be added as additional bootstrap peer. Reduces SPOF risk. \\
\rowcolor{lightgray}
Discovery hub & \prt & mDNS works on same LAN. Beyond LAN, needs protocol extension. \\
Message forwarding & \yes & As an always-on peer, can forward messages between peers via app-level relay. Especially valuable: always available as relay for any connected peer. \\
\bottomrule
\end{tabular}
\caption{Custom supernode capabilities. Gains message forwarding via app-level peer relay.}
\end{table}

\textbf{Supernode as relay hub}: An always-on supernode with a public IP is an ideal app-level relay peer. Since it is always reachable and connected to many peers, it can forward messages for any pair of connected peers. This is particularly valuable for:
\begin{itemize}
  \item Peers behind symmetric NAT where Circuit Relay + DCUtR fails
  \item Cross-protocol bridging (IPv4$\leftrightarrow$IPv6) without requiring a dual-stack Circuit Relay peer
  \item BLE-only peers that have a co-located peer with libp2p connectivity to the supernode
\end{itemize}

\textbf{Complexity}: Running a libp2p supernode requires running a full \texttt{dart\_libp2p} node with all subsystems (AutoNAT, Circuit Relay, Noise, etc.). Non-trivial to deploy and maintain.

% =============================================================================
\section{Mobility Analysis}
% =============================================================================

\subsection{M1: WiFi to Cellular Handoff}

\begin{table}[H]
\centering
\small
\begin{tabular}{L{3cm} L{2cm} C{1.5cm} L{5cm}}
\toprule
\textbf{Transition} & \textbf{Transport} & \textbf{LP2P} & \textbf{Details} \\
\midrule
WiFi (PIv4) $\to$ LTE (PIv4) & UDX/TCP & \no$^*$ & \textbf{Connection dies}. UDX and TCP are bound to the 4-tuple (src IP, src port, dst IP, dst port). When IP changes, all connections are severed. Recovery requires: reconnect to bootstrap, new relay reservation, re-establish all peer connections. \textbf{Recovery: 10--30 seconds.} \\
\rowcolor{lightgray}
WiFi (PIv4) $\to$ LTE (PuIv6) & UDX/TCP & \no$^*$ & Worse: protocol family also changes. Full reconnection required. \\
WiFi (PuIv6) $\to$ LTE (PIv4) & UDX/TCP & \no$^*$ & Downgrade from public to NATed. All connections lost; must re-establish via relay. \\
\rowcolor{lightgray}
WiFi + BLE $\to$ LTE (no BLE peer) & BLE & \yes & BLE unaffected by IP transition. \\
\bottomrule
\end{tabular}
\caption{WiFi$\to$cellular handoff. $^*$No connection migration---all connections die on IP change.}
\end{table}

\textbf{This is libp2p's biggest weakness for mobile}: Neither UDX nor TCP support connection migration. When the device's IP changes:
\begin{enumerate}
  \item All existing connections (UDX and TCP) are severed instantly
  \item The node must reconnect to the bootstrap peer from scratch
  \item AutoNAT must re-probe reachability
  \item If NATed, a new relay reservation must be made
  \item All peer connections must be re-established from scratch
  \item \textbf{Total recovery time: 10--30 seconds} of complete blackout
\end{enumerate}

\textbf{Peer relay during handoff}: If the disconnected peer has a co-located BLE peer that maintains libp2p connectivity, that BLE peer can serve as an app-level relay during the handoff period. This partially mitigates the blackout---messages can still flow through the BLE relay while libp2p reconnects.

\subsection{M2: Moving Between WiFi Networks}

\begin{table}[H]
\centering
\small
\begin{tabular}{L{3.5cm} C{1.5cm} L{5.5cm}}
\toprule
\textbf{Scenario} & \textbf{LP2P} & \textbf{Details} \\
\midrule
WiFi A $\to$ WiFi B (diff LAN) & \no$^*$ & New private IP. All connections lost. Full reconnection cycle (10--30s). \\
\rowcolor{lightgray}
Same AP, IP unchanged & \yes & No disruption. \\
WiFi A $\to$ brief gap $\to$ WiFi B & \no$^*$ & TCP/UDX timeout during gap. Even if new IP acquired quickly, connections are dead. \\
\rowcolor{lightgray}
Same LAN, different subnet & \no$^*$ & IP change = connection death. \\
\bottomrule
\end{tabular}
\caption{WiFi roaming. $^*$Connections do not survive IP changes.}
\end{table}

\subsection{M3: BLE Range Transitions}

\begin{table}[H]
\centering
\small
\begin{tabular}{L{3.5cm} C{1.5cm} C{1.5cm} L{4.5cm}}
\toprule
\textbf{Scenario} & \textbf{BLE} & \textbf{LP2P} & \textbf{Details} \\
\midrule
Peer walks into BLE range & \yes & unchanged & BLE scan discovers peer. ANNOUNCE exchanged. New transport available. \\
\rowcolor{lightgray}
Peer walks out of BLE range & \no & unchanged & BLE disconnects. LibP2P unaffected if both have IP. \\
Peer walks out, no Internet & \no & \no & Both transports lost. \textbf{Peer unreachable} unless a relay peer exists. \\
\rowcolor{lightgray}
BLE flapping (edge of range) & \dgr & unchanged & Repeated connect/disconnect. Should debounce. LibP2P provides stable fallback \textit{if IP hasn't changed}. \\
\bottomrule
\end{tabular}
\caption{BLE range transitions. BLE is independent of Internet transport.}
\end{table}

\textbf{Peer relay impact}: When BLE drops and the device is simultaneously experiencing an IP change (e.g., walking away from a WiFi AP), both direct transports fail simultaneously. However, if a third peer C maintains connectivity to both A and B (e.g., C is on the same WiFi as B and has libp2p to A's remote peer), C can relay messages during the transition.

\subsection{M4: BLE-Only Operation (No Internet)}

\begin{table}[H]
\centering
\small
\begin{tabular}{L{4cm} C{1.5cm} C{6cm}}
\toprule
\textbf{Scenario} & \textbf{Works?} & \textbf{Details} \\
\midrule
Peer next to me, both BLE-only & \yes & Full communication via BLE. Limited to BLE throughput ($\sim$1--2 KB/s). \\
\rowcolor{lightgray}
Peer 50m away, both BLE-only & \prt & At BLE range limit. Unreliable, frequent disconnects. \\
Peer 200m away, both BLE-only & \no & Out of BLE range. No alternative transport. No relay peer in between. \\
\rowcolor{lightgray}
3 peers in a line, BLE-only & \yes & Middle peer C forwards for A$\leftrightarrow$B. Single-hop relay. End-to-end encrypted. Effective range doubles. \\
Airplane mode with BLE on & \yes & BLE works independently of airplane mode radio restrictions (on most devices). \\
\rowcolor{lightgray}
Underground / no cell signal & \yes & BLE unaffected. Only local peers reachable. Relay extends range to friends-of-friends. \\
\bottomrule
\end{tabular}
\caption{BLE-only operation (no Internet). Peer relay enables the 3-peers-in-a-line scenario.}
\end{table}

\textbf{BLE-only implications}:
\begin{itemize}
  \item \textbf{Range}: $\sim$10--30m indoors, $\sim$100m outdoors per hop. With peer relay, effective range doubles to $\sim$60m indoors / $\sim$200m outdoors (single hop through a middle peer).
  \item \textbf{Throughput}: BLE GATT writes are $\sim$1--2 KB/s effective. Fine for text messages, insufficient for media. Relay adds latency but not throughput reduction (bottleneck is the BLE link itself).
  \item \textbf{Resilient}: Works in disaster scenarios, underground, in Faraday cages, on planes---anywhere Bluetooth radio works. Peer relay extends this resilience beyond direct BLE range.
\end{itemize}

% =============================================================================
\section{Combined Scenario Matrix}
% =============================================================================

\begin{longtable}{L{3.5cm} L{2.5cm} C{1cm} C{1cm} C{1cm} L{3cm}}
\toprule
\textbf{Scenario} & \textbf{Address Config} & \textbf{BLE} & \textbf{LP2P} & \textbf{Relay} & \textbf{Effective Connectivity} \\
\midrule
\endfirsthead
\toprule
\textbf{Scenario} & \textbf{Address Config} & \textbf{BLE} & \textbf{LP2P} & \textbf{Relay} & \textbf{Effective Connectivity} \\
\midrule
\endhead
\midrule
\multicolumn{6}{r}{\textit{continued on next page}} \\
\bottomrule
\endfoot
\bottomrule
\endlastfoot

% --- Co-located scenarios ---
\multicolumn{6}{l}{\textbf{Co-located Peers ($<$30m)}} \\
\midrule
Coffee shop, both on WiFi & PIv4 / PIv4 & \yes & \yes & \na & Excellent. BLE preferred. mDNS finds peer on same LAN. \\
\rowcolor{lightgray}
Coffee shop, one on LTE & PIv4 / PIv4 & \yes & \dgr & \na & BLE works. LibP2P via relay (different NATs). \\
Park, no WiFi & BLE / LTE(PIv4) & \yes & \prt & \na & BLE works. One peer has libp2p. \\
\rowcolor{lightgray}
Underground, no signal & BLE / BLE & \yes & \no & \prt$^*$ & BLE only. Relay if 3rd peer present. \\

\midrule
% --- Remote scenarios ---
\multicolumn{6}{l}{\textbf{Remote Peers (different cities)}} \\
\midrule
Both on home WiFi & PIv4 / PIv4 & \no & \dgr & \prt$^\dagger$ & Relay via bootstrap peer. Hole-punch depends on router. App relay if supernode connected. \\
\rowcolor{lightgray}
One on fiber (PuIv4) & PuIv4 / PIv4 & \no & \yes & \na & Direct connection to public peer. Fast. \\
Both on LTE & PIv4 / PIv4 & \no & \dgr & \prt$^\dagger$ & Carrier-grade NAT = symmetric. Relay-only likely. App relay through supernode if available. \\
\rowcolor{lightgray}
Both have IPv6 & PuIv6 / PuIv6 & \no & \yes & \na & Direct IPv6/UDX. Optimal for remote. \\
One offline & Any / BLE only & \no & \no & \prt$^*$ & \textbf{App relay}: if BLE peer has a co-located friend with libp2p, that friend can forward. Otherwise unreachable. \\

\midrule
% --- Mobility scenarios ---
\multicolumn{6}{l}{\textbf{Mobility}} \\
\midrule
Walking away from BLE peer, have LTE & PIv4 & $\to$\no & \dgr$^*$ & \prt & BLE drops. LibP2P \textit{if still connected}. App relay through remaining BLE peers. \\
\rowcolor{lightgray}
Walking toward peer, gaining BLE & PIv4 & $\to$\yes & \yes & \na & BLE discovered. Becomes preferred transport. \\
On train, passing through cells & PIv4 $\to$ PIv4 & \no & \no & \prt & \textbf{Constant full reconnects}. App relay through stable peer if available. \\
\rowcolor{lightgray}
Entering airplane mode + BLE & BLE & \yes & $\to$\no & \prt$^*$ & LibP2P lost. BLE works for nearby peers. Relay extends reach. \\
Leaving subway (no signal $\to$ LTE) & BLE $\to$ PIv4 & \yes & $\to$\dgr & \prt & LibP2P must bootstrap from scratch. 10--30s delay. BLE relay during gap. \\
\rowcolor{lightgray}
WiFi drops, no cellular & PIv4 $\to$ BLE & \yes & $\to$\no & \prt$^*$ & Only nearby BLE peers reachable. Relay extends to friends-of-friends. \\

\end{longtable}

\noindent $^*$Requires a co-located relay peer with BLE connectivity.\\
$^\dagger$Requires a supernode or other always-on peer connected to both parties.

% =============================================================================
\section{Gap Analysis}
% =============================================================================

\begin{enumerate}
  \item \textbf{No connection migration}: The single biggest problem. Mobile users on cellular experience constant disconnections. Train/car scenarios are effectively unusable. Peer relay through a stable BLE peer partially mitigates this.

  \item \textbf{Bootstrap SPOF}: One server down = entire transport-level relay infrastructure gone for NATed peers. Adding bootstrap peers is possible but requires configuration changes.

  \item \textbf{Cross-protocol bridging}: IPv4$\leftrightarrow$IPv6 requires a dual-stack Circuit Relay peer at the transport level. App-level peer relay can bridge this if a connected peer has connectivity to both parties (regardless of IP version).

  \item \textbf{Double identity mapping}: LibP2P PeerId $\neq$ SGN public key creates complexity and potential for identity confusion.

  \item \textbf{Relay peer availability}: App-level relay requires a Peer C connected to both A and B. In sparse networks with few peers, suitable relay peers may not exist. An always-on supernode mitigates this.

  \item \textbf{Discovery requires BLE or mDNS}: Remote peers that have never been co-located and are not on the same LAN cannot discover each other without out-of-band exchange.

  \item \textbf{Single relay hop limit}: Only one relay hop is allowed. In scenarios where A can reach C and C can reach D and D can reach B, but no single peer connects A to B, communication is impossible. This is a deliberate simplicity trade-off.
\end{enumerate}

% =============================================================================
\section{Recommendations}
% =============================================================================

\begin{table}[H]
\centering
\begin{tabular}{C{0.5cm} L{4cm} L{3cm} L{4cm}}
\toprule
\textbf{\#} & \textbf{Gap} & \textbf{Recommendation} & \textbf{Impact} \\
\midrule
1 & No connection migration & Evaluate QUIC-based transport (e.g., iroh) & Survive IP changes on mobile \\
\rowcolor{lightgray}
2 & Bootstrap SPOF & Deploy redundant bootstrap peers & Resilience for NATed peers \\
3 & Double identity mapping & Unify PeerId with SGN public key & Simpler architecture \\
\rowcolor{lightgray}
4 & Discovery beyond BLE & Address exchange via QR code, link & Enable remote-first friendships \\
5 & Relay peer availability & Deploy always-on supernode as relay hub & Guarantee relay availability \\
\rowcolor{lightgray}
6 & Peer relay is single-hop only & Evaluate multi-hop if needed for mesh scenarios & Extended range at cost of complexity \\
\bottomrule
\end{tabular}
\caption{Recommendations for identified connectivity gaps.}
\end{table}

% =============================================================================
\section{Summary}
% =============================================================================

\begin{table}[H]
\centering
\begin{tabular}{L{5cm} C{2cm} C{2cm} C{2cm}}
\toprule
\textbf{Scenario Class} & \textbf{BLE} & \textbf{LibP2P} & \textbf{App Relay} \\
\midrule
Co-located, same LAN & \yes & \yes{} (mDNS) & \na \\
Co-located, BLE range & \yes & \na & \na \\
Co-located, beyond BLE, within relay range & \no & \na & \yes \\
Remote, both NATed & \no & \dgr & \prt$^\dagger$ \\
Remote, one public & \no & \yes & \na \\
Remote, BLE-only peer & \no & \no & \yes$^*$ \\
Corporate firewall (TCP only) & \no & \yes & \na \\
Mobility (IP change) & unchanged & \no{} (reconnect) & \prt$^*$ \\
Mobility (BLE transition) & varies & unchanged & \prt$^*$ \\
Bootstrap failure & \na & \no{} (SPOF) & unchanged \\
BLE-only (no Internet) & \yes & \no & \yes$^*$ \\
\bottomrule
\end{tabular}
\caption{High-level comparison by scenario class. $^*$Requires co-located relay peer. $^\dagger$Requires supernode.}
\end{table}

The three connectivity mechanisms (BLE, libp2p, and application-level peer relay) are \textbf{complementary}:

\begin{itemize}
  \item \textbf{BLE} provides resilient local connectivity that is immune to network configuration.
  \item \textbf{LibP2P} provides global reach that is immune to physical proximity.
  \item \textbf{Peer relay} bridges the gap between the two: it extends BLE range, bridges transport boundaries (BLE$\leftrightarrow$libp2p), and provides a fallback path when direct transports fail.
\end{itemize}

The primary remaining gaps are:
\begin{itemize}
  \item \textbf{No connection migration}: LibP2P's UDX/TCP connections die on IP change. This is the single biggest problem for mobile users. A QUIC-based transport would solve this.
  \item \textbf{Bootstrap SPOF}: The bootstrap peer is a single point of failure for NATed peers' transport-level relay.
  \item \textbf{No store-and-forward}: By design. Relay is real-time only. If the recipient is unreachable via any path (direct or relayed), the message fails.
\end{itemize}

Peer relay mitigates some of these gaps (e.g., a co-located BLE peer can relay during a libp2p handoff), but it does not eliminate the fundamental limitations of the transport layer. The relay is only as good as the connectivity of the relay peer itself.

\end{document}
